\documentclass[11pt]{article}
\usepackage[a4paper, tmargin=2.5cm, lmargin=2.8cm, rmargin=2.8cm, bmargin=2.5cm]{geometry}
\usepackage{hyperref}

\hypersetup{
    colorlinks=true,
    linkcolor=black,
    filecolor=magenta,
    urlcolor=blue,
    citecolor=black,
}

\usepackage{graphicx}
\usepackage{amsmath}
\usepackage{amssymb}
\usepackage{array}
\usepackage[fontset=fandol]{ctex}
\usepackage{enumitem}
\usepackage{booktabs}
\usepackage{caption}
\usepackage{xcolor}
\pagestyle{plain}

\begin{document}

\begin{center}
{\LARGE\bfseries Prof-Finder：基于LLM与语义匹配的导师推荐系统} \\[6pt]
{\large 技术选型与产品体系结构设计报告}
\end{center}

\vspace{5mm}

%==========================================================
\section{技术选型}
%==========================================================

\subsection{选型原则}

技术选型遵循以下原则：

\begin{enumerate}
    \item \textbf{轻量化优先}：选择轻量级框架与嵌入式数据库，降低部署与运维复杂度。
    \item \textbf{Python 生态优先}：核心业务逻辑集中在 Python 生态，保证机器学习、大模型等库的兼容性。
    \item \textbf{现代前端工程化}：前端采用 TypeScript + 组合式 API + 原子化 CSS。
    \item \textbf{异步优先}：后端全链路支持 async/await，避免 I/O 密集型操作阻塞请求。
    \item \textbf{跨平台兼容}：开发和使用都同时支持到 macOS 与 Windows，照顾到组内的所有成员。
\end{enumerate}

\begin{table}[htbp]
  \centering
  \caption{系统开发环境（已在 macOS 26 Tahoe 完成环境搭建）}
  \begin{tabular}{ll}
    \toprule
    开发环境 & 环境参数 \\
    \midrule
    开发操作系统 & 已在 macOS 26 Tahoe 完成环境搭建 \\
    内存 & 16GB \\
    CPU & M2 (ARM64) \\
    Python 环境管理 & Conda（prof-finder 环境） \\
    Python 版本 & 3.11.15 \\
    Python 包管理器 & Poetry 2.3.4 \\
    后端框架 & FastAPI 0.109.2 \\
    ASGI 服务器 & Uvicorn 0.27.1 \\
    ORM 框架 & SQLAlchemy 2.0.46 \\
    数据库 & SQLite \\
    任务队列 & Huey 3.0.0 \\
    语义嵌入框架 & sentence-transformers 5.2.3 \\
    语义嵌入模型 & allenai-specter \\
    LLM SDK & OpenAI SDK 1.109.1 \\
    爬虫库 & scholarly 1.7.11 / requests 2.32.5 / BeautifulSoup4 4.14.3 \\
    PDF 解析 & pymupdf4llm 0.3.4 \\
    前端框架 & Vue 3.5.24 \\
    TypeScript 版本 & 5.9.3 \\
    构建工具 & Vite 7.3.1 \\
    UI 组件库 & Naive UI 2.43.2 \\
    CSS 框架 & Tailwind CSS 4.2.4 \\
    状态管理 & Pinia 3.0.4 \\
    HTTP 客户端 & Axios 1.13.3 \\
    实时通信 & EventSource（SSE）/ sse-starlette 2.4.1 \\
    Node.js 版本 & v24.15.0 \\
    npm 版本 & 11.12.1 \\
    IDE & VSCode / PyCharm \\
    版本控制 & Git 2.46.1 \\
    代码格式化 & Black 23.12.1 \\
    代码检查 & Flake8 6.1.0 / MyPy 1.19.1 \\
    测试框架 & pytest 7.4.4 \\
    \bottomrule
  \end{tabular}
  \label{tab:dev_env_macos}
\end{table}

\begin{table}[htbp]
  \centering
  \caption{系统开发环境（已在 Windows 11 完成环境搭建）}
  \begin{tabular}{ll}
    \toprule
    开发环境 & 环境参数 \\
    \midrule
    开发操作系统 & 已在 Windows 11 完成环境搭建 \\
    内存 & 16GB \\
    CPU & Intel i7-12700H \\
    Python 环境管理 & Conda（prof-finder 环境） \\
    Python 版本 & 3.11.15 \\
    Python 包管理器 & Poetry 2.3.4 \\
    后端框架 & FastAPI 0.109.2 \\
    ASGI 服务器 & Uvicorn 0.27.1 \\
    ORM 框架 & SQLAlchemy 2.0.46 \\
    数据库 & SQLite \\
    任务队列 & Huey 3.0.0 \\
    语义嵌入框架 & sentence-transformers 5.2.3 \\
    语义嵌入模型 & allenai-specter \\
    LLM SDK & OpenAI SDK 1.109.1 \\
    爬虫库 & scholarly 1.7.11 / requests 2.32.5 / BeautifulSoup4 4.14.3 \\
    PDF 解析 & pymupdf4llm 0.3.4 \\
    前端框架 & Vue 3.5.24 \\
    TypeScript 版本 & 5.9.3 \\
    构建工具 & Vite 7.3.1 \\
    UI 组件库 & Naive UI 2.43.2 \\
    CSS 框架 & Tailwind CSS 4.2.4 \\
    状态管理 & Pinia 3.0.4 \\
    HTTP 客户端 & Axios 1.13.3 \\
    实时通信 & EventSource（SSE）/ sse-starlette 2.4.1 \\
    Node.js 版本 & v24.15.0 \\
    npm 版本 & 11.12.1 \\
    IDE & VS Code / PyCharm \\
    版本控制 & Git 2.46.1 \\
    代码格式化 & Black 23.12.1 \\
    代码检查 & Flake8 6.1.0 / MyPy 1.19.1 \\
    测试框架 & pytest 7.4.4 \\
    \bottomrule
  \end{tabular}
  \label{tab:dev_env_windows}
\end{table}

\subsection{技术选型说明}

\subsubsection{后端}

\textbf{FastAPI} 作为 Web 框架，天然支持 async/await，适合处理 I/O 密集型操作（LLM 调用、爬虫请求、文件解析）；内置 Pydantic 数据验证与 OpenAPI 文档自动生成。

\textbf{SQLAlchemy 2.0 + SQLite} 作为持久层。SQLAlchemy 提供成熟 ORM；SQLite 免去独立数据库服务的运维开销。后期可无缝迁移至 PostgreSQL。

\textbf{sentence-transformers} 加载 allenai-specter 模型（基于 SciBERT 微调），将教授研究方向与学生简历映射至同一语义空间，计算余弦相似度。

\textbf{OpenAI SDK} 调用 DeepSeek API（OpenAI 兼容接口），完成简历解析、论文摘要与邮件生成。

\textbf{Huey} 作为轻量级任务队列，配合 SQLite 实现后台任务的可靠调度。

\textbf{scholarly + BeautifulSoup4 + requests} 构成爬虫工具链：scholarly 负责 Google Scholar，BeautifulSoup4 负责大学院系网页解析。

\textbf{pymupdf4llm} 将 PDF 论文转换为 Markdown，供 LLM 进行摘要与关键词提取。

\subsubsection{前端}

\textbf{Vue 3 + Composition API + TypeScript} 为核心框架，提供灵活的逻辑复用与静态类型检查。

\textbf{Vite} 利用原生 ES 模块实现极速冷启动与 HMR。

\textbf{Naive UI} 提供 90+ 组件，天然 Tree Shaking 与 TypeScript 支持。

\textbf{Tailwind CSS} 以原子化方式构建样式，减少手写 CSS。

\textbf{Pinia} 作为 Vuex 继任者管理全局状态，提供简洁 API 与 TypeScript 推导。

\textbf{SSE} 接收后端异步任务进度推送，比 WebSocket 更轻量，符合单向推送场景。

% %==========================================================
% \section{产品体系结构设计}
% %==========================================================

% \subsection{整体架构}

% 系统采用\textbf{前后端分离}的四层架构：

% \begin{center}
% \small
% \begin{tabular}{lll}
% \toprule
% \textbf{层次} & \textbf{职责} & \textbf{关键技术} \\
% \midrule
% 表示层 & 用户界面、路由、交互 & Vue 3 + Naive UI + Tailwind \\
% 服务层 & API 网关、认证、业务编排 & FastAPI + JWT + SSE \\
% 领域层 & 爬虫、解析、匹配、LLM 调用 & scholarly + sentence-transformers + OpenAI SDK \\
% 持久层 & 数据存储、缓存、迁移 & SQLAlchemy + SQLite \\
% \bottomrule
% \end{tabular}
% \end{center}

% 前端构建产物为纯静态文件，由 Nginx 托管；后端作为独立进程通过 Uvicorn 运行。
% Nginx 同时充当反向代理，将 \texttt{/api/*} 请求转发至后端，其余请求指向静态资源。

% \subsection{前端架构}

% 前端采用页面—组件两级结构，Vue Router（history 模式）管理路由，Pinia 管理全局状态（认证、任务、设置）。
% 路由守卫在 \texttt{router.beforeEach} 中依次校验认证状态、强制修改密码标志与管理员权限。
% 异步任务的实时进度通过 SSE（EventSource）获取，TaskPanel 组件统一展示并支持页面刷新恢复。

% \subsection{后端架构}

% 后端按 API 层、业务层、数据层、配置层组织。
% 全部路由以 \texttt{/api} 为前缀，按资源划分为认证、简历、教授、匹配、邮件、任务、文档输入、用户设置八个路由组。

% 认证采用 JWT 双令牌（Access Token 15 分钟 + Refresh Token 7 天），密码 Bcrypt 哈希存储，FastAPI 依赖注入实现声明式权限校验。

% 耗时操作（爬取、匹配、LLM 调用）通过异步任务系统处理：API 立即返回 \texttt{task\_id}，后台通过 \texttt{asyncio} 执行，SSE 实时推送进度，支持任务取消。

% \subsection{数据架构}

% 核心实体包括：User（用户）、UserSettings（设置）、UserProfile（简历）、Professor（教授）、MatchRecord（匹配记录）、SourceInput（文档输入）。
% 所有业务实体携带 \texttt{user\_id}，数据访问层强制过滤，实现行级多用户数据隔离。

% 采用轻量级增量迁移方案：启动时通过 \texttt{PRAGMA table\_info} 检测缺失列并自动补全，后续可迁移至 Alembic。
% 教授语义向量持久化至数据库，信息变更时自动置为无效并强制重算。

% \subsection{核心模块}

% \textbf{爬虫模块}采用注册表模式：Google Scholar 通过 scholarly 库获取，大学院系通过 BeautifulSoup4 爬取，新增爬虫只需实现基类并注册，前端动态获取可用列表。

% \textbf{语义匹配}通过 allenai-specter 将教授与简历文本分别编码为语义向量，计算余弦相似度，输出排名。教授向量首次计算后缓存复用。

% \textbf{LLM 集成}覆盖简历解析、论文摘要、邮件生成三类任务。LLM 不可用时自动降级为启发式方法。API Key 支持用户级与全局级两级配置。

%==========================================================
\section{产品体系结构设计}
%==========================================================

\subsection{整体架构}

系统采用\textbf{前后端分离}的经典分层架构，共分为四层：

\begin{center}
\small
\begin{tabular}{lp{4.5cm}p{5cm}}
\toprule
\textbf{层次} & \textbf{职责} & \textbf{关键技术} \\
\midrule
表示层（Presentation） & 用户界面、交互逻辑、路由 & Vue 3 + Naive UI + Tailwind \\
服务层（Application） & API 网关、认证、业务编排 & FastAPI + JWT + SSE \\
领域层（Domain） & 爬虫、简历解析、语义匹配、LLM 调用 & scholarly + sentence-transformers + OpenAI SDK \\
持久层（Persistence） & 数据存储、缓存、迁移 & SQLAlchemy + SQLite \\
\bottomrule
\end{tabular}
\end{center}

前后端通过 HTTP/HTTPS 协议通信，API 统一以 \texttt{/api} 为前缀。
前端构建产物为纯静态文件（HTML/CSS/JS），可由任意 Web 服务器（Nginx、Caddy）或 CDN 托管。
后端作为独立进程运行，通过 ASGI 服务器（Uvicorn）对外暴露接口。

\begin{center}
\begin{tabular}{ccccc}
\toprule
\textbf{浏览器} & $\rightarrow$ & \textbf{Nginx} & $\rightarrow$ & \textbf{静态资源（dist/）} \\
\midrule
& & & $\rightarrow$ & \textbf{API 服务（Uvicorn:8000）} \\
\bottomrule
\end{tabular}
\end{center}

\subsection{前端体系结构}

\subsubsection{路由设计}

采用 Vue Router 的 history 模式，路由映射如下：

\begin{center}
\begin{tabular}{lll}
\toprule
\textbf{路由路径} & \textbf{页面} & \textbf{权限} \\
\midrule
\texttt{/login} & 登录 & 公开 \\
\texttt{/register} & 注册 & 公开 \\
\texttt{/change-password} & 修改密码 & 已登录（首登强制跳转） \\
\texttt{/profile} & 简历管理 & 已登录 \\
\texttt{/professor} & 教授管理 & 已登录 \\
\texttt{/professor/:id/edit} & 教授信息编辑 & 已登录 \\
\texttt{/match} & 匹配结果 & 已登录 \\
\texttt{/letter} & 邮件管理 & 已登录 \\
\texttt{/settings} & 用户设置 & 已登录 \\
\texttt{/admin/users} & 用户管理 & 仅管理员 \\
\bottomrule
\end{tabular}
\end{center}

\subsubsection{实时通信}

异步任务（爬取、匹配、LLM 调用）通过 SSE 实现进度推送：

\begin{enumerate}
    \item 用户在前端发起操作（如"爬取教授"）
    \item 前端调用 API 获取 \texttt{task\_id}，立即返回
    \item 前端创建 \texttt{EventSource} 连接到 \texttt{/api/tasks/\{task\_id\}/progress}
    \item 后端异步执行任务，通过 SSE 推送进度事件（progress / complete / failed）
    \item 前端 TaskPanel 组件实时更新进度条与状态
\end{enumerate}

\subsection{后端体系结构}

\subsubsection{分层设计}

后端采用\textbf{四层架构}组织代码：

\begin{center}
\small
\begin{tabular}{lp{4.5cm}p{4cm}}
\toprule
\textbf{层次} & \textbf{目录} & \textbf{职责} \\
\midrule
API 层 & \texttt{api/} & 路由定义、请求验证、依赖注入、SSE 端点 \\
业务层 & \texttt{crawler/, matcher/, llm/, parser/} & 爬虫调度、语义匹配、LLM 调用、简历解析 \\
数据层 & \texttt{db/, models/} & 数据库连接、ORM 模型、CRUD 操作 \\
配置层 & \texttt{prompts/} & YAML Prompt 模板、配置常量 \\
\bottomrule
\end{tabular}
\end{center}

\subsubsection{API 路由规划}

所有路由统一以 \texttt{/api} 为前缀，按资源类型划分：

\begin{center}
\begin{tabular}{lll}
\toprule
\textbf{路由前缀} & \textbf{资源} & \textbf{核心端点} \\
\midrule
\texttt{/api/auth/} & 认证 & login, register, refresh, me \\
\texttt{/api/profiles/} & 简历 & 上传、解析预览、CRUD、激活 \\
\texttt{/api/professors/} & 教授 & Scholar 爬取、大学批量爬取、CRUD、信息编辑 \\
\texttt{/api/match/} & 匹配 & 运行匹配、获取结果 \\
\texttt{/api/letters/} & 邮件 & 生成、查看、编辑 \\
\texttt{/api/tasks/} & 异步任务 & SSE 进度、取消、批量清理 \\
\texttt{/api/source-inputs/} & 文档输入 & PDF 上传、ArXiv 提交、重试 \\
\texttt{/api/settings/} & 用户设置 & 读取、更新 API Key 等 \\
\bottomrule
\end{tabular}
\end{center}

\subsubsection{异步任务系统}

系统中爬取、LLM 调用、匹配计算均为耗时操作，设计异步任务系统如下：

\begin{itemize}
    \item \textbf{任务管理器}：维护全局任务状态字典，每个任务包含唯一 ID、类型、进度、状态与事件队列
    \item \textbf{任务创建}：API 端点接收请求后，立即创建 \texttt{TaskState} 并返回 \texttt{task\_id}
    \item \textbf{后台执行}：通过 \texttt{asyncio.create\_task} 启动协程，阻塞 I/O（如 scholarly 请求）通过 \texttt{asyncio.to\_thread} 卸载到线程池
    \item \textbf{进度推送}：每个任务维护一个 \texttt{asyncio.Queue}，SSE 端点从中读取事件并推送到前端
    \item \textbf{取消支持}：任务循环中检查 \texttt{cancel\_requested} 标志，支持用户主动取消
\end{itemize}

支持的任务类型包括：单教授爬取、批量爬取、大学院系爬取、语义匹配、单封邮件生成、批量邮件生成、论文摘要。

\subsection{数据体系结构}

\subsubsection{核心数据模型}

我们设计了六类核心实体，其关系如图~\ref{fig:er} 所示。

\begin{table}[htbp]
\centering
\small
\begin{tabular}{p{2.2cm}p{5cm}p{5.3cm}}
\toprule
\textbf{实体} & \textbf{核心字段} & \textbf{关联} \\
\midrule
User & id, username, password\_hash, is\_admin & $\rightarrow$ UserSettings, UserProfile, Professor \\
UserSettings & user\_id, deepseek\_api\_key, request\_delay & $\rightarrow$ User（一对一） \\
UserProfile & user\_id, title, name, education, skills, is\_active & $\rightarrow$ User（多对一） \\
& & $\rightarrow$ MatchRecord \\
Professor & user\_id, name, affiliation, research\_interests, & $\rightarrow$ User（多对一） \\
& publications, embedding, h\_index & $\rightarrow$ MatchRecord, SourceInput \\
MatchRecord & user\_profile\_id, professor\_id, score, letter\_content & $\rightarrow$ UserProfile, Professor \\
SourceInput & source\_type, canonical\_id, extracted\_markdown, status & $\rightarrow$ Professor（可选） \\
\bottomrule
\end{tabular}
\caption{核心数据模型及其关联}
\label{fig:er}
\end{table}

\subsubsection{爬虫模块设计}

爬虫模块采用\textbf{注册表模式}设计，支持可扩展的多数据源接入：

\begin{center}
\small
\begin{tabular}{lp{4.5cm}p{3.5cm}}
\toprule
\textbf{数据源} & \textbf{技术方案} & \textbf{输出} \\
\midrule
Google Scholar & scholarly 库，按 ID 搜索并填充基本信息、论文列表与引用指标 & Professor 结构化字典 \\
大学院系网站 & BeautifulSoup4 + requests，继承 UniversityCrawlerBase 基类 & Professor 列表字典 \\
\bottomrule
\end{tabular}
\end{center}

新增大学爬虫只需：实现 \texttt{UniversityCrawlerBase} 子类 $\rightarrow$ 在注册表中添加一行映射。
前端通过 \texttt{GET /api/professors/university-crawlers} 动态获取可用爬虫列表，实现\textbf{开闭原则}（对扩展开放，对修改关闭）。

%==========================================================
\section{小组成员开发环境构建情况}
%==========================================================

全体成员已按照第~\ref{tab:dev_env_macos}、\ref{tab:dev_env_windows} 节的配置完成开发环境搭建。
小组分工与负责方向如下：

\begin{center}
\begin{tabular}{lll}
\toprule
\textbf{姓名} & \textbf{角色} & \textbf{开发环境搭建情况} \\
\midrule
郝飞洋 & 项目发起人 & 已在 macOS 26 Tahoe 完成环境搭建 \\
李晨璞 & 项目经理 & 已在 Windows 11 完成环境搭建 \\
谢寒迪 & 后端架构师 & 已在 Windows 11 完成环境搭建 \\
李仁昌 & 数据库工程师 & 已在 macOS 26 Tahoe 完成环境搭建 \\
白雨欣 & 前端开发负责人 & 已在 Windows 11 完成环境搭建 \\
彭子飏 & AI 与算法工程师 & 已在 Windows 11 完成环境搭建 \\
李昊文 & 数据工程师 & 已在 Windows 11 完成环境搭建 \\
银睿迪 & 测试 \& 文档 \& 部署 & 已在 Windows 11 完成环境搭建 \\
\bottomrule
\end{tabular}
\end{center}

各成员开发环境均已验证通过，项目具备进入下一阶段的条件。

\end{document}
